% !PRINTED 31
\weberpage{31}
Sind aber $\alpha$ und $\beta$ in zwei verschiedenen Gruppen, etwa in $Q_\nu$ und $Q'_\nu$ enthalten, so schliessen wir so: Weil $\beta$ in $P_\nu$ enthalten und $Q_\nu$ ein Normaltheiler von $P_\nu$ ist, so ist auch
\weq{11}{\beta^{-1}\alpha\beta = \alpha'}
in $Q_\nu$ enthalten. Daraus folgt:
\weq{12}{\beta\alpha'\alpha^{-1} = \alpha\beta\alpha^{-1} = \beta',}
und weil $\alpha$ in $P_\nu$ enthalten und $Q'_\nu$ ein Normaltheiler von $P_\nu$ ist, so ist auch $\beta'$ in $Q'_\nu$ enthalten. Aus (12) aber folgt
\[ \alpha'\alpha^{-1} = \beta^{-1}\beta' = \delta, \]
und demnach ist das Element $\delta$ sowohl in $Q_\nu$ als in $Q'_\nu$ enthalten. Es kann also nach dem, was unter 1) bewiesen ist, nur das Einheitselement sein, d. h. es ist
\[ \alpha' = \alpha, \quad \beta' = \beta, \]
und demnach folgt aus (11), dass $\alpha\beta = \beta\alpha$ sein muss, wie bewiesen werden sollte.

\begin{center}
\rule{3cm}{0.4pt}
\end{center}
% [The source prints this plain horizontal rule below the paragraph, marking the close of \S.\,8; the remainder of the page is blank.]
